% GENERATED FILE. Edit canonical lessons, then run npm run generate:books.
% canonical-source-hash: fnv1a64:1cb7e690f3620c9a
% canonical-lessons: ZH-C265-haixiu, ZH-C265-yonggan, ZH-C265-dafang, ZH-C265-xiaoqi, ZH-C265-jinzhang, ZH-R265-review-260-265

\chapter{Shy or Brave}
\label{ch:zh-shy-or-brave}
\hlchaptermodality{265}

\begin{chapteropening}
\textbf{By the end of this chapter:} \emph{I can say that someone is shy, brave, generous, stingy or nervous.}

Complete the last lesson of chapter 265: I can say that someone is shy, brave, generous, stingy or nervous.
\end{chapteropening}

\section[hàixiū]{hàixiū — shy; bashful}

\label{lesson:ZH-C265-haixiu}



\begin{quote}
Before the new one: say the Mandarin for afraid that, then the Mandarin for as a matter of fact.
\end{quote}

\subsection*{You'll want to know: hàixiū}
\emph{hàixiū} — \textquotedblleft{}shy; bashful\textquotedblright{}.

\begin{quote}
\emph{nàge háizi hěn hàixiū} — That child is shy.
\end{quote}

Its characters come later; for now it is a word for the ear.

\begin{cousinweb}
\emph{hàixiū} is built from \emph{hài} 'to harm' and \emph{xiū} 'shame', both inherited from Middle Chinese, \emph{hajH} and \emph{sjuw} (Guangyun, in Baxter's transcription), which glosses the first 'to harm' and the second 'shame'.
\end{cousinweb}

\begin{grammarlens}[title={What you've built}]
One more word for describing things.
\end{grammarlens}

\subsection*{Your turn}
\begin{itemize}
\raggedright
  \item \emph{Say it:} \emph{hàixiū}
  \item \emph{Say it:} \emph{hàixiū}, once more
  \item \emph{Say it:} \emph{nàge háizi hěn hàixiū} — That child is shy.
  \item \emph{Recall:} say the Mandarin for afraid that, then the Mandarin for as a matter of fact, then say \emph{hàixiū} again
\end{itemize}

\begin{tcolorbox}[breakable,colback=teal!4,colframe=teal!35!black,title={Before you move on}]
What does hàixiū mean? (\textquotedblleft{}Shy; bashful\textquotedblright{}.) Say it once more.
\end{tcolorbox}
\section[yǒnggǎn]{yǒnggǎn — brave}

\label{lesson:ZH-C265-yonggan}



\begin{quote}
Before the new one: say the Mandarin for as a matter of fact, then the Mandarin for shy.
\end{quote}

\subsection*{You'll want to know: yǒnggǎn}
\emph{yǒnggǎn} — \textquotedblleft{}brave\textquotedblright{}.

\begin{quote}
\emph{tā hěn yǒnggǎn} — He's very brave.
\end{quote}

Its characters come later; for now it is a word for the ear.

Two third tones meet in this word: yǒng is said with a rising tone before gǎn, and the spelling keeps both marks.

\begin{cousinweb}
\emph{yǒnggǎn} is built from \emph{yǒng} and \emph{gǎn}, both inherited from Middle Chinese, \emph{yowngX} and \emph{kamX} (Guangyun, in Baxter's transcription), which glosses the first 'fierce' and the second 'brave' and 'to dare'.
\end{cousinweb}

\begin{grammarlens}[title={What you've built}]
One more word for describing things.
\end{grammarlens}

\subsection*{Your turn}
\begin{itemize}
\raggedright
  \item \emph{Say it:} \emph{yǒnggǎn}
  \item \emph{Say it:} \emph{yǒnggǎn}, once more
  \item \emph{Say it:} \emph{tā hěn yǒnggǎn} — He's very brave.
  \item \emph{Recall:} say the Mandarin for as a matter of fact, then the Mandarin for shy, then say \emph{yǒnggǎn} again
\end{itemize}

\begin{tcolorbox}[breakable,colback=teal!4,colframe=teal!35!black,title={Before you move on}]
What does yǒnggǎn mean? (\textquotedblleft{}Brave\textquotedblright{}.) Say it once more.
\end{tcolorbox}
\section[dàfang]{dàfang — generous; open-handed}

\label{lesson:ZH-C265-dafang}



\begin{quote}
Before the new one: say the Mandarin for shy, then the Mandarin for brave.
\end{quote}

\subsection*{You'll want to know: dàfang}
\emph{dàfang} — \textquotedblleft{}generous; open-handed\textquotedblright{}.

\begin{quote}
\emph{wǒ gēge hěn dàfang} — My brother is very generous.
\end{quote}

Its characters come later; for now it is a word for the ear.

\begin{cousinweb}
\emph{dàfang} is built from \zh{大} 'big' and \emph{fāng} 'square', both inherited from Middle Chinese, \emph{daH} and \emph{pjang} (Guangyun, in Baxter's transcription), which glosses the second 'the four quarters' and 'upright'. In the word the second loses its tone; how the pair came to mean generous is not settled.
\end{cousinweb}

\begin{grammarlens}[title={What you've built}]
One more word for describing things.
\end{grammarlens}

\subsection*{Your turn}
\begin{itemize}
\raggedright
  \item \emph{Say it:} \emph{dàfang}
  \item \emph{Say it:} \emph{dàfang}, once more
  \item \emph{Say it:} \emph{wǒ gēge hěn dàfang} — My brother is very generous.
  \item \emph{Recall:} say the Mandarin for shy, then the Mandarin for brave, then say \emph{dàfang} again
\end{itemize}

\begin{tcolorbox}[breakable,colback=teal!4,colframe=teal!35!black,title={Before you move on}]
What does dàfang mean? (\textquotedblleft{}Generous; open-handed\textquotedblright{}.) Say it once more.
\end{tcolorbox}
\section[\texorpdfstring{\zh{小气} (xiǎoqì)}{小气 (xiǎoqì)}]{\zh{小气} (xiǎoqì) — stingy; mean (with money)}

\label{lesson:ZH-C265-xiaoqi}



\begin{quote}
Before the new one: say the Mandarin for brave, then the Mandarin for generous.
\end{quote}

\subsection*{You'll want to know: \zh{小气}}
\textbf{\zh{小气}} — \emph{xiǎoqì} — \textquotedblleft{}stingy; mean (with money)\textquotedblright{}.

\begin{quote}
\textbf{\zh{他很小气}} \emph{tā hěn xiǎoqì} — He's stingy.
\end{quote}

\begin{cousinweb}
\textbf{\zh{小气}} (\emph{xiǎoqì}) is built from \zh{小} 'small' and \zh{气} 'breath', both inherited from Middle Chinese, \emph{sjewX} and \emph{khj+jH} (Guangyun, in Baxter's transcription), which glosses the second 'breath'.
\end{cousinweb}

\begin{grammarlens}[title={What you've built}]
One more word for describing things.
\end{grammarlens}

\subsection*{Your turn}
\begin{itemize}
\raggedright
  \item \emph{Say it:} \emph{xiǎoqì}
  \item \emph{Say it:} \emph{xiǎoqì}, once more
  \item \emph{Say it:} \emph{tā hěn xiǎoqì} — He's stingy.
  \item \emph{Recall:} say the Mandarin for brave, then the Mandarin for generous, then say \emph{xiǎoqì} again
\end{itemize}

\begin{tcolorbox}[breakable,colback=teal!4,colframe=teal!35!black,title={Before you move on}]
What does \zh{小气} mean? (\textquotedblleft{}Stingy; mean (with money)\textquotedblright{}.) Say it once more.
\end{tcolorbox}
\section[jǐnzhāng]{jǐnzhāng — nervous; tense}

\label{lesson:ZH-C265-jinzhang}



\begin{quote}
Before the new one: say the Mandarin for generous, then the Mandarin for stingy.
\end{quote}

\subsection*{You'll want to know: jǐnzhāng}
\emph{jǐnzhāng} — \textquotedblleft{}nervous; tense\textquotedblright{}.

\begin{quote}
\emph{míngtiān kǎoshì wǒ hěn jǐnzhāng} — I have an exam tomorrow, and I'm nervous.
\end{quote}

Its characters come later; for now it is a word for the ear.

\begin{cousinweb}
\emph{jǐnzhāng} is built from \emph{jǐn} 'tight' and \emph{zhāng} 'to stretch', both inherited from Middle Chinese, \emph{kjinX} and \emph{trjang} (Guangyun, in Baxter's transcription), which glosses the first 'tightly bound' and the second 'to stretch out'.
\end{cousinweb}

\begin{grammarlens}[title={What you've built}]
One more word for describing things.
\end{grammarlens}

\subsection*{Your turn}
\begin{itemize}
\raggedright
  \item \emph{Say it:} \emph{jǐnzhāng}
  \item \emph{Say it:} \emph{jǐnzhāng}, once more
  \item \emph{Say it:} \emph{míngtiān kǎoshì wǒ hěn jǐnzhāng} — I have an exam tomorrow, and I'm nervous.
  \item \emph{Recall:} say the Mandarin for generous, then the Mandarin for stingy, then say \emph{jǐnzhāng} again
\end{itemize}

\begin{tcolorbox}[breakable,colback=teal!4,colframe=teal!35!black,title={Before you move on}]
What does jǐnzhāng mean? (\textquotedblleft{}Nervous; tense\textquotedblright{}.) Say it once more.
\end{tcolorbox}
\section[(dialogue)]{Looking back: from What I Can Do to Shy or Brave}

\label{lesson:ZH-R265-review-260-265}



\begin{quote}
Say the Mandarin for ability and what someone is capable of.
\end{quote}

\subsection*{Your turn}
\begin{itemize}
\raggedright
  \item \emph{Say it:} ability, what someone is capable of, an expert, technique, to make progress
  \item \emph{Say it:} know-how, at one time, to study abroad, overseas, to experience for yourself
  \item \emph{Say it:} tennis, table tennis, an opponent, to defeat, to cheer someone on
  \item \emph{Say it:} a manager, to be in charge, to ask for time off, sick leave, annual leave
  \item \emph{Say it:} to turn down, a shame, sorry to say, afraid that, as a matter of fact
  \item \emph{Say it:} shy, brave, generous, stingy, nervous
  \item \emph{Read it:} \zh{假} — each one aloud, then a word that holds it
\end{itemize}

\begin{tcolorbox}[breakable,colback=teal!4,colframe=teal!35!black,title={Before you move on}]
Ability? (\emph{nénglì}.) What someone is capable of? (\textbf{\zh{本事}}, \emph{běnshi}.) An expert? (\textbf{\zh{高手}}, \emph{gāoshǒu}.) Technique? (\emph{jìshù}.) To make progress? (\emph{jìnbù}.) Know-how? (\emph{jīngyàn}.) At one time? (\emph{céngjīng}.) To study abroad? (\emph{liúxué}.) Overseas? (\textbf{\zh{国外}}, \emph{guówài}.) To experience for yourself? (\emph{tǐyàn}.) Tennis? (\emph{wǎngqiú}.) Table tennis? (\emph{pīngpāngqiú}.) An opponent? (\textbf{\zh{对手}}, \emph{duìshǒu}.) To defeat? (\emph{dǎbài}.) To cheer someone on? (\emph{jiāyóu}.) A manager? (\emph{jīnglǐ}.) To be in charge? (\emph{fùzé}.) To ask for time off? (\textbf{\zh{请假}}, \emph{qǐngjià}.) Sick leave? (\textbf{\zh{病假}}, \emph{bìngjià}.) Annual leave? (\textbf{\zh{年假}}, \emph{niánjià}.) To turn down? (\emph{jùjué}.) A shame? (\emph{kěxī}.) Sorry to say? (\emph{bàoqiàn}.) Afraid that? (\emph{kǒngpà}.) As a matter of fact? (\emph{qíshí}.) Shy? (\emph{hàixiū}.) Brave? (\emph{yǒnggǎn}.) Generous? (\emph{dàfang}.) Stingy? (\textbf{\zh{小气}}, \emph{xiǎoqì}.) Nervous? (\emph{jǐnzhāng}.)
\end{tcolorbox}
